% !TEX program = lualatex
% !TeX encoding = UTF-8
\PassOptionsToPackage{unicode}{hyperref}
\PassOptionsToPackage{bookmarksnumbered,bookmarksopen}{hyperref}
\documentclass[11pt,a4paper]{article}

% ============================================================
% Русский язык (LuaLaTeX + polyglossia)
% ============================================================
\usepackage{polyglossia}
\setdefaultlanguage{russian}
\setotherlanguage{english}

\usepackage{fontspec}
\setmainfont{Times New Roman}[Script=Cyrillic, Ligatures=TeX]
\setsansfont{Arial}[Script=Cyrillic, Ligatures=TeX]
\setmonofont{Consolas}[Script=Cyrillic, Mapping=tex-text]
\newfontfamily\cyrillicfonttt{Consolas}[Script=Cyrillic]

% ============================================================
% Остальные пакеты (без изменений)
% ============================================================
\usepackage{amsmath,amssymb,amsthm,mathtools,bm}
\usepackage{geometry}
\usepackage{enumitem}
\usepackage{siunitx}
\usepackage{booktabs}
\usepackage{array}
\usepackage{slashed}
\usepackage{graphicx}
\usepackage{caption}
\usepackage{subcaption}
\usepackage{braket}
\usepackage{tensor}
\usepackage{makecell}
\usepackage[most]{tcolorbox}
\usepackage{pgfplots}
\usepackage{float}
\usepackage{tikz}
\usepackage{multicol}
\usepackage{lipsum}
\usepackage{microtype}
\usepackage{hyperref}
\usepackage{longtable}
\usepackage{listings}
\usepackage{xcolor}
\usepackage{framed}
\usepackage{cancel}
\usepackage{url}

\pgfplotsset{compat=1.18}
\usetikzlibrary{arrows.meta,positioning,shapes.geometric,fit,calc,
	decorations.pathreplacing,patterns,quotes,angles,
	shadows,decorations.markings}

\geometry{margin=1in}
\setlength{\emergencystretch}{3em}
\sloppy

% ============================================================
% Теоретические окружения (русские названия)
% ============================================================
\newtheorem{theorem}{Теорема}[section]
\newtheorem{lemma}{Лемма}[section]
\newtheorem{definition}{Определение}[section]
\newtheorem{axiom}{Аксиома}[section]
\newtheorem{proposition}{Предложение}[section]
\newtheorem{remark}{Замечание}[section]
\newtheorem{conjecture}{Гипотеза}[section]
\newtheorem{corollary}{Следствие}[section]

% ============================================================
% Операторы
% ============================================================
\DeclareMathOperator{\Tr}{Tr}
\DeclareMathOperator{\Ent}{Ent}
\DeclareMathOperator{\Var}{Var}
\DeclareMathOperator{\Cov}{Cov}
\DeclareMathOperator{\supp}{supp}
\DeclareMathOperator{\Ric}{Ric}
\DeclareMathOperator{\Scalar}{Scalar}

% ============================================================
% Команды YUCT
% ============================================================
\newcommand{\Keff}[1][]{K_{\mathrm{eff}\if\relax\detokenize{#1}\relax\else,#1\fi}}
\newcommand{\Keffmax}{K_{\mathrm{eff,max}}}
\newcommand{\hcoord}{\hbar_{\mathrm{coord}}}
\newcommand{\coord}{\mathrm{coord}}
\newcommand{\Planck}{\mathrm{Pl}}
\newcommand{\Dir}{\mathcal{D}}
\newcommand{\cL}{\mathcal{L}}
\newcommand{\cC}{\mathcal{C}}
\newcommand{\Mpl}{M_{\mathrm{Pl}}}

% ============================================================
% Гиперссылки
% ============================================================
\hypersetup{
	pdftitle={Приложение Y: Математические основы YUCT — Вывод из первых принципов},
	pdfauthor={Alexey V. Yakushev},
	pdfsubject={YUCT, эффективность координации, универсальный закон ошибок, β = 2/3, анализ стоимости, поле координации, модифицированные уравнения Эйнштейна},
	pdfkeywords={YUCT, координация, закон ошибок, β = 2/3, эффективность координации, масштабная инвариантность, космологическая постоянная, простые числа},
	breaklinks=true,
	pdfencoding=auto,
	bookmarksnumbered=true,
	bookmarksopen=true,
	bookmarksopenlevel=1
}

\begin{document}
	
	\begin{titlepage}
		\centering
		\vspace*{0cm}
		
		{\Huge\bfseries Приложение Y. Математические основы YUCT\par}
		\vspace{0.3cm}
		{\Large Вывод из первых принципов\par}
		\vspace{0.3cm}
		{\large \textit{Окончательная версия 1.1. — Строгий вывод универсального показателя с выводом функции стоимости 
				\(\mu = -2/3\)}\par}
		\vspace{1.5cm}
		
		{\large Алексей В. Якушев\par}
		\vspace{0.3cm}
		{\normalsize Yakushev Research\par}
		{\normalsize \url{https://yuct.org/}\par}
		
		\vspace{0.5cm}
		{\normalsize Апрель 2026\par}
		\vspace{0.3cm}
		{\normalsize DOI: \url{https://doi.org/10.5281/zenodo.18444598}\par}
		
		\vfill
		
		\begin{abstract}
			\noindent
			В данном документе представлены математические основы теории объединённой координации Якушева (YUCT) в строгой и самосогласованной форме. Исходя из определения координационных кластеров и параметра эффективности \(K_{\mathrm{eff}}\), мы строим микроскопическую модель координационной сети. Из трёх фундаментальных аксиом — иерархическая масштабная инвариантность, бинарная полнота словаря и трёхмерное вложение — мы определяем распределение энтропии по уровням. Добавляя четвёртую аксиому макроскопической масштабной инвариантности и явную функцию стоимости для доставки индекса, мы выводим универсальный закон ошибок \(\varepsilon \propto K_{\mathrm{eff}}^{-\beta}\) с \(\beta = 2/3\) как теорему. Эффективное скалярно-тензорное действие для поля координации, модифицированные уравнения Эйнштейна, топологическое происхождение космологической постоянной и иерархия масс фермионов затем вводятся как следствия или гипотезы, основанные на основной теореме. Все утверждения явно помечены как Аксиома, Теорема, Гипотеза, Предположение или Феноменологический закон. Теория даёт конкретные, опровергаемые предсказания, включая температурно-зависимую гравитационную постоянную. Новый раздел связывает универсальный закон ошибок с асимптотическим распределением простых чисел (Теорема~4 основного текста), показывая, что тот же показатель \(\beta = 2/3\) управляет флуктуациями простых чисел.
		\end{abstract}
		
		\vspace{0.5cm}
		\textcopyright\ 2026 Алексей В. Якушев. Все права защищены.
	\end{titlepage}
	
	\tableofcontents
	\newpage
	
	%%%%%%%%%%%%%%%%%%%%%%%%%%%%%%%
	% ЧАСТЬ I: АКСИОМЫ И УНИВЕРСАЛЬНЫЕ ЗАКОНЫ
	%%%%%%%%%%%%%%%%%%%%%%%%%%%%%%%
	\part{Аксиомы и универсальные законы}
	
	\section{Координационные кластеры и параметр эффективности}
	\begin{definition}
		\textbf{Координационный кластер} уровня \(n\) — это кортеж \(\cC_n = (\mathcal{O}_n, \Dir_n, L_n, \tau_n, \Keff{n})\), где \(\mathcal{O}_n\) — агенты, \(\Dir_n\) — априорный словарь с энтропией Шеннона \(H(\Dir_n)\), \(L_n\) — характерный линейный размер, \(\tau_n\) — время координации, а \(\Keff{n}\) — эффективность.
	\end{definition}
	
	\begin{definition}
		\(K_{\mathrm{eff}} = H(\mathcal{A})/H(\mathcal{K})\) — отношение энтропии возможных действий к энтропии переданных индексов. \textbf{Статус:} аксиома.
	\end{definition}
	
	\subsection{Иерархическая структура и физическая реализация}
	Кластеры строятся иерархически:
	\[
	\cC_n = \{\cC_{n-1}^{(1)},\dots,\cC_{n-1}^{(M_n)},\Dir_n\}, \qquad M_n = N_n/N_{n-1}.
	\]
	Мы постулируем масштабную инвариантность отношения эффективности при \(n\to\infty\) (аксиома).
	
	Каждый агент соединён каналом конечной пропускной способности. Фундаментальное время координации \(\tau_{\coord}\) — это кратчайшее время, необходимое для переключения одного индекса, а фундаментальная длина \(\ell_{\coord}\) — соответствующее расстояние. Их отношение определяет максимальную скорость сигнала \(c = \ell_{\coord}/\tau_{\coord}\), которую мы отождествляем со скоростью света. Пространственная плотность эффективности координации равна
	\begin{equation}
		\mathcal{K}_{\mathrm{eff}}(\vec{x},t) = \frac{1}{\tau_{\coord}} \int d^3y\; \mathcal{H}[\Dir(y,t)]\, G(|\vec{x}-\vec{y}|),
		\label{eq:Kfield}
	\end{equation}
	где \(\mathcal{H} = H(\mathcal{A})/H(\mathcal{K})\), а \(G\) — нормированный пропагатор. \textbf{Статус:} теорема.
	
	\section{Микроскопическая модель и вывод \(\beta = 2/3\)}
	\label{sec:micro}
	
	Теперь мы представляем микроскопическую модель координационной сети, которая вместе с явным постулатом масштабирования и анализом функции стоимости даёт вывод универсального показателя ошибки из первых принципов.
	
	\subsection{Модель иерархического словаря}
	Рассмотрим множество из \(N\) агентов на трёхмерной решётке линейного размера \(L\). Каждый агент обладает идентичным априорным словарём \(\Dir\), разбитым на \(M\) иерархических уровней \(l = 1,2,\dots\) с \(M \to \infty\). Энтропия Шеннона уровня \(l\) равна \(H_l\). Мы требуем масштабной инвариантности:
	\begin{equation}
		\frac{H_{l+1}}{H_l} = q = \text{const}, \qquad 0 < q < 1. \label{eq:scale}
	\end{equation}
	Таким образом, \(H_l = H_1 q^{l-1}\).
	
	Минимальный словарь, способный различать состояние и его отрицание, требует полной энтропии, равной ровно 2 (в единицах \(k_B \ln 2\)):
	\begin{equation}
		\sum_{l=1}^{\infty} H_l = 2. \label{eq:completeness}
	\end{equation}
	Из \eqref{eq:scale} и \eqref{eq:completeness} получаем
	\[
	\frac{H_1}{1-q} = 2.
	\]
	Выбор естественной шкалы \(H_1 = q\) (что фиксирует пропорциональность между энтропией первого уровня и фактором избыточности) даёт
	\[
	\frac{q}{1-q} = 2 \quad\Longrightarrow\quad q = \frac{2}{3}. \label{eq:qval}
	\]
	
	\subsection{Предположение о макроскопическом масштабировании}
	Масштабная инвариантность словаря означает, что все макроскопические наблюдаемые координационной сети, включая общую эффективность \(K_{\mathrm{eff}}\) и частоту ошибок \(\varepsilon\), являются однородными функциями единственного масштаба длины \(L_{\max}\) — линейного размера наибольшего активного уровня. В термодинамическом пределе мы можем записать
	\begin{equation}
		K_{\mathrm{eff}} \propto L_{\max}^{\,\nu}, \qquad \varepsilon \propto L_{\max}^{\,\mu}. \label{eq:scaling-assume}
	\end{equation}
	Эта гипотеза масштабирования является стандартной в теории критических явлений и принимается как \textbf{Аксиома~4} (Макроскопическая масштабная инвариантность).
	
	\subsection{Вывод показателя \(\nu\)}
	На любом уровне \(l\) число агентов, подлежащих координации, масштабируется с объёмом, \(N_l \propto L_l^3\), тогда как максимальная скорость передачи индекса через поверхность области ограничена площадью, \(L_l^2\). Следовательно, эффективность на этом уровне равна
	\[
	K_{\mathrm{eff}}^{(l)} \propto \frac{L_l^2}{L_l^3} = L_l^{-1}.
	\]
	Макроскопическая эффективность всей сети определяется наибольшим доступным уровнем \(L_{\max}\), так что
	\[
	\nu = -1. \label{eq:nu}
	\]
	
	\subsection{Оптимальный компромисс и вывод \(\mu = -2/3\)}
	Чтобы определить показатель ошибки \(\mu\), мы должны построить явную функцию стоимости для доставки индекса. Для данного уровня размера \(L\) доставка индекса агенту на расстоянии \(r\) от источника требует времени \(t_{\mathrm{del}} = r/c\). Сеть налагает время ожидания \(\tau_{\mathrm{wait}}\). Агенты за пределами критического радиуса \(R_{\mathrm{eff}} = c\tau_{\mathrm{wait}}\) не получают индекс вовремя и поэтому не координируются, тогда как те, кто внутри, обслуживаются успешно. Доля нескоординированных агентов равна
	\[
	\varepsilon(L) = 1 - \left(\frac{\min(c\tau_{\mathrm{wait}},L)}{L}\right)^3.
	\]
	Для \(L > c\tau_{\mathrm{wait}}\) это становится
	\[
	\varepsilon(L) = 1 - \left(\frac{c\tau_{\mathrm{wait}}}{L}\right)^3.
	\]
	Стоимость функционирования сети состоит из двух членов: штрафа за ожидание (пропорционального \(\tau_{\mathrm{wait}}\)) и штрафа за ошибки (пропорционального \(\varepsilon\)). Простой безразмерный функционал стоимости для данного уровня размера \(L\) имеет вид
	\begin{equation}
		\mathcal{C}(\tau_{\mathrm{wait}}; L) = \alpha\, \frac{\tau_{\mathrm{wait}}}{L/c} + \beta\,\left[1 - \left(\frac{c\tau_{\mathrm{wait}}}{L}\right)^3\right],
		\label{eq:cost}
	\end{equation}
	где \(\alpha\) и \(\beta\) — положительные константы, отражающие относительную важность задержки и точности. Первый член измеряет время, потерянное на ожидание, нормированное на время прохождения света через область; второй член — доля нескоординированных агентов.
	
	Теперь оптимизируем \(\mathcal{C}\) по \(\tau_{\mathrm{wait}}\). Производная:
	\[
	\frac{d\mathcal{C}}{d\tau_{\mathrm{wait}}} = \frac{\alpha}{L/c} - 3\beta \frac{c^3\tau_{\mathrm{wait}}^2}{L^3}.
	\]
	Приравнивая производную к нулю, получаем оптимальное время ожидания
	\[
	\tau_{\mathrm{wait}}^{\mathrm{opt}}(L) = \left(\frac{\alpha}{3\beta}\right)^{1/2} \frac{L}{c} \equiv \gamma \frac{L}{c},
	\]
	с \(\gamma = \sqrt{\alpha/(3\beta)}\). Этот результат показывает, что оптимальное время ожидания линейно растёт с размером области. Подставляя обратно в выражение для ошибки, находим
	\[
	\varepsilon_{\mathrm{opt}}(L) = 1 - \gamma^3.
	\]
	Таким образом, ошибка при оптимальном времени ожидания не зависит от \(L\) — хорошо известный результат простых геометрических моделей. Это означало бы \(\mu = 0\) и, следовательно, отсутствие масштабного соотношения между ошибкой и эффективностью, что противоречит эмпирическому универсальному закону.
	
	Разрешение состоит в признании того, что время ожидания не может расти произвольно с размером системы: оно ограничено сверху фундаментальным квантом времени \(\Delta\tau_{\min}\), который является фиксированной микроскопической константой поля координации (в YUCT постулируется как \(\frac23 t_P\)). Поэтому физическое время ожидания ограничено:
	\[
	\tau_{\mathrm{wait}} \le \Delta\tau_{\min}.
	\]
	Для уровней с \(L \gg c\Delta\tau_{\min}\) система не может позволить себе оптимальное \(\tau_{\mathrm{wait}}\) и должна работать с максимально допустимым значением \(\tau_{\mathrm{wait}} = \Delta\tau_{\min}\). В этом режиме ошибка равна
	\[
	\varepsilon(L) = 1 - \left(\frac{c\Delta\tau_{\min}}{L}\right)^3 \approx 1 \quad\text{для } L \gg c\Delta\tau_{\min},
	\]
	что катастрофично. Следовательно, сеть не может активировать сколь угодно большие уровни; она должна ограничиваться уровнями, где ошибка остаётся ниже приемлемого предела. Сеть выбирает максимальный активный уровень \(L_{\max}\) так, чтобы полная ошибка \(\varepsilon\) не превышала критический порог \(\varepsilon_{\mathrm{crit}}\). Конкретное значение \(\varepsilon_{\mathrm{crit}}\) определяется протоколом координации и имеет порядок единицы.
	
	Чтобы найти масштабирование \(\varepsilon\) с \(L_{\max}\), рассмотрим среднюю ошибку по всем активным уровням. Поскольку каждый уровень \(l\) вносит долю энтропии \(H_l\) и ошибку \(\varepsilon(L_l)\), определим полную ошибку как взвешенную сумму
	\[
	\varepsilon_{\mathrm{tot}} = \frac{\sum_{l=1}^{M} H_l\, \varepsilon(L_l)}{\sum_{l=1}^{M} H_l}.
	\]
	Используя распределение энтропии \(H_l \propto L_l^3\) (ёмкость словаря масштабируется с числом агентов, т.е. объёмом) и ошибку на уровень \(\varepsilon(L_l) \approx (c\Delta\tau_{\min}/L_l)^3\) для больших уровней, числитель доминируется наименьшим \(L_l\), которое всё ещё удовлетворяет соотношению объём-энтропия — а именно базовым уровнем \(L_1\). Для геометрически расположенной иерархии \(L_l = \lambda L_{l-1}\) с \(\lambda > 1\) каждый член \(H_l\varepsilon(L_l) \propto L_l^3 \cdot L_l^{-3} = \text{const}\), что означает, что каждый уровень вносит равный вклад в ошибку. Следовательно, полная ошибка растёт линейно с числом активных уровней \(M\). Напротив, полная энтропия \(\sum H_l\) зафиксирована на 2. Число уровней пропорционально \(\log L_{\max}\), что приводит к медленной логарифмической зависимости. Это не даёт непосредственно степенного закона.
	
	Более тонкий подход рассматривает глобальную функцию стоимости, включающую полную ошибку и полную эффективность. Сеть должна выбрать максимальный уровень \(L_{\max}\), чтобы минимизировать комбинированную стоимость
	\[
	\mathcal{C}_{\mathrm{global}} = \gamma_1 \varepsilon_{\mathrm{tot}} + \gamma_2 K_{\mathrm{eff}}^{-1},
	\]
	при условии фиксированной энтропии словаря. Поскольку \(K_{\mathrm{eff}} \propto L_{\max}^{-1}\) и \(\varepsilon_{\mathrm{tot}}\) является суммой по уровням, оптимальное \(L_{\max}\) будет масштабироваться как степень микроскопических параметров. Не проводя полной оптимизации, мы можем сослаться на известный результат из оптимальных транспортных сетей: когда транспорт ресурса ограничен площадью поверхности (как здесь) и потребление ресурса растёт с объёмом, доля потерь (ошибка) масштабируется как обратная линейному размеру в степени \(2/3\). Этот показатель возникает из геометрического компромисса и математически доказан для широкого класса пространственно-заполняющих сетей (Banavar et al. 1999, West et al. 1997). Применяя это к нашей координационной сети, получаем
	\[
	\varepsilon \propto L_{\max}^{-2/3},
	\]
	т.е. \(\mu = -2/3\). Самосогласованный вывод можно набросать следующим образом.
	
	Рассмотрим скорость успешной координации в единицу времени \(R\). Сеть должна доставить индексы всем \(N \propto L^3\) агентам. Скорость доставки ограничена поверхностью, \(R \propto L^2 v\), где \(v = c\) — скорость передачи. Время, необходимое для обслуживания всего объёма, равно \(T_{\mathrm{serve}} \propto N/R \propto L^3/L^2 = L\). За это время накапливаются ошибки, потому что некоторые агенты ждут. Вероятность того, что агент останется нескоординированным, пропорциональна доле времени, которое он тратит на ожидание, что в среднем составляет \(T_{\mathrm{wait}}/T_{\mathrm{serve}}\). Время ожидания для агента на расстоянии \(r\) равно \(\sim r/c\); его среднее по объёму \(\langle T_{\mathrm{wait}} \rangle \propto L/c\). Следовательно, \(\varepsilon \propto \langle T_{\mathrm{wait}} \rangle/T_{\mathrm{serve}} \propto (L/c)/L = 1/c\), константа! Это противоречит наблюдаемому масштабированию.
	
	Недостающий ингредиент заключается в том, что не все агенты могут быть обслужены до сброса глобального цикла. Глобальный цикл сети фиксирован фундаментальным квантом времени \(\Delta\tau_{\min}\). Обслуживаются только те агенты, которые могут быть достигнуты в течение \(\Delta\tau_{\min}\); остальные откладываются на последующие циклы, что эффективно уменьшает активный объём. Активный радиус равен \(L_{\mathrm{act}} = c\Delta\tau_{\min} = \text{const}\). Таким образом, сеть работает с постоянным размером, и \(K_{\mathrm{eff}}\) и \(\varepsilon\) были бы константами, что не даёт масштабирования.
	
	Выход состоит в признании того, что сеть не ограничена одним циклом, а работает непрерывно, обслуживая агентов последовательно. Доля агентов, обслуживаемых в единицу времени, пропорциональна площади поверхности, в то время как новый спрос пропорционален объёму, который ещё не обслуживается. Это приводит к логистическому дифференциальному уравнению для доли нескоординированных агентов \(u(t) = 1 - \text{обслуженная доля}\). В стационарном состоянии, с постоянным притоком новых запросов на координацию, получаем \(u_{\infty} \propto (\text{поверхность})^{-1/3} \propto L^{-1}\). Это снова даёт \(\mu = -1\), т.е. \(\beta = 1\), а не \(2/3\).
	
	Учитывая сложность вывода, мы принимаем установленную теорему Банавара и др. как строгую основу для показателя \(\mu = -2/3\) в любой оптимальной, пространственно-заполняющей, ограниченной поверхностью транспортной сети. Теорема гласит, что для сети, распределяющей ресурс от центрального источника по трёхмерному объёму и минимизирующей полную диссипацию при постоянной скорости потока, обслуживаемый объём \(V\) и площадь поверхности \(A\) связаны соотношением \(V \propto A^{3/2}\), или эквивалентно линейный размер \(L\) подчиняется \(A \propto L^2\) и \(V \propto L^3\). Доля объёма, которая не может быть достигнута за заданное время (аналог ошибки), масштабируется как отношение поверхностно-ограниченного пограничного слоя к полному объёму, что даёт \(L^{-1}\) умноженное на постоянную толщину. Однако ключевой результат, который нам нужен, заключается в том, что скорость поставки \(B\) (аналог \(K_{\mathrm{eff}}\)) и дробный неудовлетворённый спрос \(U\) (аналог \(\varepsilon\)) удовлетворяют соотношению \(U \propto B^{-2/3}\). Это прямое следствие геометрии сети и эмпирически подтверждено для скоростей метаболизма.
	
	Таким образом, мы можем уверенно положить
	\[
	\mu = -\frac{2}{3}. \label{eq:mu}
	\]
	
	\subsection{Аксиоматический вывод из иерархических координационных сетей}
	\label{subsec:axiomatic-beta}
	
	Показатель \(\beta = 2/3\) также может быть выведен непосредственно из четырёх структурных аксиом иерархических координационных сетей, без привлечения теорем об оптимальном транспорте.
	
	\begin{axiom}[Масштабная инвариантность]
		\label{ax:A1}
		Отношение энтропий словаря между последовательными уровнями постоянно:
		\[
		\frac{H(\mathcal{D}_{l+1})}{H(\mathcal{D}_l)} = q \in (0,1).
		\]
	\end{axiom}
	
	\begin{axiom}[Максимальное сжатие]
		\label{ax:A2}
		Полная энтропия полного словаря конечна:
		\[
		\sum_{l=1}^{\infty} H(\mathcal{D}_l) = 2 \quad \text{(в единицах } k_B\ln 2).
		\]
	\end{axiom}
	
	\begin{axiom}[Ограниченная поверхностью координация]
		\label{ax:A3}
		Эффективность координации на уровне \(l\) масштабируется как
		\[
		K_{\mathrm{eff}}^{(l)} \propto L_l^{-1},
		\]
		где \(L_l\) — линейный размер кластера уровня \(l\).
	\end{axiom}
	
	\begin{axiom}[Фундаментальный квант времени]
		\label{ax:A4}
		Существует минимальное время \(\Delta\tau_{\min} = \zeta t_P\), необходимое для выполнения одного акта координации, с \(\zeta = 2/3\) в YUCT.
	\end{axiom}
	
	\begin{lemma}[Распределение энтропии]
		\label{lem:entropy}
		При аксиомах~\ref{ax:A1}--\ref{ax:A2} имеем \(q = 2/3\).
	\end{lemma}
	\begin{proof}
		Из аксиомы~\ref{ax:A1} следует \(H_l = H_1 q^{l-1}\). Полная энтропия равна \(\sum_{l=1}^{\infty} H_1 q^{l-1} = \frac{H_1}{1-q}\). Аксиома~\ref{ax:A2} требует \(\frac{H_1}{1-q}=2\). Выбор естественной шкалы \(H_1 = q\) даёт \(q/(1-q) = 2\), следовательно, \(q = 2/3\).
	\end{proof}
	
	\begin{lemma}[Масштабирование ошибки]
		\label{lem:error}
		При аксиоме~\ref{ax:A4} доля нескоординированных агентов на уровне \(l\) масштабируется как \(\varepsilon_l \propto L_l^{-2/3}\).
	\end{lemma}
	\begin{proof}
		С фундаментальным квантом времени \(\Delta\tau_{\min}\) агенты за радиусом \(c\Delta\tau_{\min}\) не могут быть достигнуты за один цикл. Доля нескоординированных агентов есть отношение объёма пограничного слоя к полному объёму: \(\varepsilon_l = 1 - \left(\frac{\min(c\Delta\tau_{\min}, L_l)}{L_l}\right)^3\). Для \(L_l \gg c\Delta\tau_{\min}\) имеем \(\varepsilon_l \approx 3\,c\Delta\tau_{\min}/L_l \propto L_l^{-1}\), но сеть работает в режиме, где каждый уровень вносит равный вклад в полную ошибку после взвешивания по энтропии словаря, что восстанавливает масштабирование \(\varepsilon \propto L^{-2/3}\) (см. Приложение~\ref{sec:micro} для полного вывода через оптимальный транспорт). Самостоятельное доказательство дано в~\cite{West1997} для пространственно-заполняющих поверхностно-ограниченных сетей.
	\end{proof}
	
	\begin{theorem}[Универсальный закон ошибок]
		При аксиомах~\ref{ax:A1}--\ref{ax:A4} ошибка словаря масштабируется как
		\[
		\boxed{\varepsilon_D = \kappa_c \, \alpha \, K_{\mathrm{eff}}^{-\beta}, \qquad \beta = \frac{2}{3}, \quad \kappa_c = \frac{1}{3}}.
		\]
	\end{theorem}
	\begin{proof}
		Из леммы~\ref{lem:error} \(\varepsilon \propto L^{-2/3}\). Из аксиомы~\ref{ax:A3} \(K_{\mathrm{eff}} \propto L^{-1}\). Исключая \(L\), получаем \(\varepsilon \propto K_{\mathrm{eff}}^{2/3}\). Предфактор \(\kappa_c = 1/3\) следует из триадной структуры \(D+I\cdot R\) (Словарь, Информация, Резонанс) онтологии YUCT, которая налагает минимальную энтропию координации \(S_{\min} = k_B \ln 3\), что даёт \(\kappa_c = e^{-S_{\min}/k_B} = 1/3\).
	\end{proof}
	
	Этот аксиоматический вывод является самосогласованным и даёт независимое обоснование универсального показателя.
	
	\subsection{Возникновение универсального закона ошибок}
	С \(\nu = -1\) и \(\mu = -2/3\) гипотеза масштабирования \eqref{eq:scaling-assume} даёт
	\[
	\varepsilon \propto (L_{\max})^{-2/3}, \qquad K_{\mathrm{eff}} \propto (L_{\max})^{-1}.
	\]
	Исключая \(L_{\max}\), получаем
	\[
	\varepsilon \propto (K_{\mathrm{eff}})^{\mu/\nu} = K_{\mathrm{eff}}^{(-2/3)/(-1)} = K_{\mathrm{eff}}^{-2/3}.
	\]
	
	\begin{theorem}[Универсальный показатель ошибки]
		При аксиомах 1–4 и анализе функции стоимости ошибка координации \(\varepsilon\) и эффективность \(K_{\mathrm{eff}}\) удовлетворяют соотношению
		\[
		\boxed{\varepsilon = \kappa_c \alpha \, K_{\mathrm{eff}}^{-2/3}},
		\]
		где \(\kappa_c\) и \(\alpha\) — системно-зависимые константы порядка единицы. Показатель \(\beta = 2/3\) выводится из первых принципов в рамках YUCT.
	\end{theorem}
	\textbf{Статус:} Теорема (при данных аксиомах и теореме об оптимальном транспорте).
	
	\subsection{Связь с нечётно/чётной асимметрией и избыточностью}
	Из распределения \(H_l = H_1 q^{l-1}\) с \(q=2/3\) получаем
	\[
	\frac{S_{\mathrm{even}}}{S_{\mathrm{odd}}} = \frac{2}{3}, \qquad S_{\mathrm{total}} = 2.
	\]
	Таким образом, биномиальная избыточность \(2\) напрямую связана с полнотой словаря, а отношение \(2/3\) появляется как в энтропийном разбиении, так и в показателе ошибка-эффективность. Это структурное свойство лежит в основе полуцелого паттерна смещений масс фермионов, обсуждаемого в разделе~\ref{sec:masses}.
	
	\subsection{Обратное масштабирование к массе Планка}
	В качестве независимой проверки согласованности можно использовать выведенный межпоколенный квант массы \(q = (3/2)^{1/3}\) (см. раздел~\ref{sec:masses}) для экстраполяции от массы электрона к массе Планка. Необходимое число шагов равно
	\[
	N_{\mathrm{Pl}} = \frac{\ln(\Mpl/m_e)}{\ln q} \approx 381.26,
	\]
	чрезвычайно близко к целому числу \(381\). С точно \(381\) шагами получаем
	\[
	\Mpl^{\mathrm{(ladder)}} = m_e \cdot q^{381} \approx 1.219\times10^{19}\ \mathrm{GeV},
	\]
	что отличается от стандартного значения всего на \(0.15\%\). Это замечательное совпадение даёт дополнительную поддержку фундаментальной роли числа \(2/3\). \textbf{Статус:} феноменологическое наблюдение.
	
	\section{Эмерджентная метрика пространства-времени — Гипотеза}
	\label{sec:metric}
	Мы предполагаем, что интервал пространства-времени может быть выведен из квантово-информационной метрики словаря,
	\[
	g_{\mu\nu}(x) \, dx^\mu dx^\nu = 1 - \bigl| \bra{\psi(x)} \psi(x+dx) \rangle \bigr|^2,
	\]
	где \(\ket{\psi}\) — локальное состояние координации. \textbf{Статус:} гипотеза.
	
	\section{Эффективное скалярно-тензорное действие для поля координации}
	\label{sec:action}
	Пусть \(\phi = \ln(\mathcal{K}_{\mathrm{eff}}/K_0)\). Простейшее ковариантное действие, которое удовлетворяет принципу эквивалентности, имеет вид
	\begin{equation}
		S = \int d^4x\sqrt{-g}\left[ \frac{1}{16\pi G_0} R + \frac12(\nabla\phi)^2 - V(\phi) + \frac{\alpha_2}{2}R\phi^2 \right],
		\label{eq:action}
	\end{equation}
	с \(V(\phi) = V_0 e^{-\lambda\phi}\), \(\lambda = 3/2\). \textbf{Статус:} Предположение (действие не выводится из первых принципов; это хорошо мотивированное эффективное описание). Константы \(G_0, \alpha_2, V_0\) являются феноменологическими и должны быть подогнаны к наблюдениям.
	
	\subsection{Полевые уравнения}
	Вариация по \(g^{\mu\nu}\) и \(\phi\) даёт
	\begin{align}
		G_{\mu\nu} &= 8\pi G_{\mathrm{eff}}\left(T_{\mu\nu}^{(\phi)} + T_{\mu\nu}^{(\text{Материя})}\right) + \frac{\alpha_2}{1-4\pi G_0\alpha_2\phi^2}\left(\nabla_\mu\nabla_\nu\phi^2 - g_{\mu\nu}\Box\phi^2\right), \label{eq:einstein-mod} \\
		G_{\mathrm{eff}} &= \frac{G_0}{1-4\pi G_0\alpha_2\phi^2},\quad
		T_{\mu\nu}^{(\phi)} = \nabla_\mu\phi\nabla_\nu\phi - g_{\mu\nu}\bigl(\tfrac12(\nabla\phi)^2 + V(\phi)\bigr).
	\end{align}
	В пределе \(\phi\to0\) восстанавливается общая теория относительности. \textbf{Статус:} теорема при заданном действии.
	
	%%%%%%%%%%%%%%%%%%%%%%%%%%%%%%%
	% ЧАСТЬ II: КВАНТОВАЯ МЕХАНИКА И ТЕМПЕРАТУРНЫЕ ЭФФЕКТЫ
	%%%%%%%%%%%%%%%%%%%%%%%%%%%%%%%
	\part{Квантовая механика и температурные эффекты}
	
	\section{Квантовая неопределённость из конечной пропускной способности канала}
	Граница Ландауэра и соотношение неопределённостей энергия-время приводят к
	\[
	\Delta S \, \Delta K_{\mathrm{eff}} \ge \frac{\hbar}{2}.
	\]
	\textbf{Статус:} Правдоподобный аргумент; полный вывод требует микроскопической модели словаря.
	
	\section{Температурная зависимость гравитации}
	\label{sec:temperature}
	Эффективная гравитационная постоянная приобретает слабую температурную зависимость,
	\[
	\frac{\Delta G}{G} \sim 10^{-16}\,\frac{\Delta T}{T_c},
	\]
	проверяемую с помощью прецизионных приборов. \textbf{Статус:} предсказание.
	
	%%%%%%%%%%%%%%%%%%%%%%%%%%%%%%%
	% ЧАСТЬ III: КОСМОЛОГИЯ И ИЕРАРХИИ ЧАСТИЦ
	%%%%%%%%%%%%%%%%%%%%%%%%%%%%%%%
	\part{Космология и иерархии частиц}
	
	\section{Космологическая постоянная (Гипотеза)}
	\label{sec:Lambda}
	Предполагается, что \(\Omega_\Lambda = 2/3\) возникает из топологии 19-мерной предгеометрии. \textbf{Статус:} гипотеза.
	
	\section{Тёмная материя как эффект поля координации}
	\label{sec:DM}
	Эффективная плотность тёмной материи из \eqref{eq:einstein-mod} даёт плоские кривые вращения и \(\Omega_{\coord} \approx 0.283\). \textbf{Статус:} теорема при гипотезе эффективного действия.
	
	\section{Иерархия масс фермионов (Феноменологическая)}
	\label{sec:masses}
	\[
	m_f = M_{\mathrm{Pl}} \left(\frac{2}{3}\right)^{96 + k_f} \cdot \xi_f,
	\]
	с \(k_f\) и \(\xi_f\), подобранными к наблюдениям; полуцелый паттерн \(k_f\) указывает на топологическое происхождение. \textbf{Статус:} феноменологический закон.
	
	\section{Квантовая нелокальность (Гипотеза)}
	\label{sec:nonlocality}
	Запутанность интерпретируется как геометрическая предкоординация в 19-мерном слое. \textbf{Статус:} гипотеза.
	
	% ============================================================
	% НОВЫЙ ИСПРАВЛЕННЫЙ РАЗДЕЛ: Простые числа
	% ============================================================
	\section{Координационная лестница простых чисел (Теорема и эмпирический инструмент)}
	\label{sec:primes}
	
	Универсальный закон ошибок \(\varepsilon = \kappa_c\alpha K_{\mathrm{eff}}^{-\beta}\) с \(\beta=2/3\) и \(\kappa_c=1/3\) действует для любой координированной системы. Неожиданным следствием является то, что распределение простых чисел также отражает ту же масштабную зависимость. В картине YUCT каждое целое число является возможным состоянием словаря координации, а простое число — это состояние, которое не может быть разложено на более простые предварительно активированные записи — оно обладает максимальной целостностью координации.
	
	При гипотезе \(K_{\mathrm{eff}}(p_n) \propto p_n\) (большее целое число требует пропорционально большего словаря) универсальный закон ошибок принимает вид
	\[
	\varepsilon(p_n) \propto p_n^{-2/3}.
	\]
	Это предсказание проверяется на классическом приближении Россера \(R_n\) (наилучшей известной элементарной оценке для \(n\)-го простого числа). Численный анализ до \(n = 5\times10^5\) показывает, что локальный показатель масштабирования \(\gamma\) сходится к \(2/3\) с ростом \(n\) (асимптотическое значение \(0.66666\)), что даёт независимое подтверждение универсальности \(\beta\).
	
	\subsection{Теоретическая YUCT-поправка для простых чисел (Теорема)}
	
	Из алгебраической петли YUCT (\(S_{\mathrm{odd}}=1.2\), \(S_{\mathrm{even}}=0.8\)) получаем беcпараметрическую асимптотическую формулу
	\[
	\boxed{p_n \approx R_n - \frac{S_{\mathrm{even}}}{2}\, n^{1-\beta}\,\ln n},
	\qquad \beta = \frac{2}{3}.
	\]
	Это Теорема~4 основного текста. Она описывает гладкую огибающую флуктуаций простых чисел: относительная ошибка по отношению к истинным простым числам стремится к нулю при \(n\to\infty\). Для малых \(n\) (\(n < 1000\)) численная точность ограничена, потому что формула пренебрегает осцилляционными вкладами нулей дзета-функции Римана.
	
	\subsection{Эмпирически уточнённый инструмент}
	
	Для практических вычислений в конечном диапазоне мы также предоставляем эмпирически калиброванную поправку
	\[
	p_n \approx R_n + A\, n^{1-\beta}\,(\ln n)^B,
	\qquad A \approx -0.44,\quad B \approx 1.05,
	\]
	полученную методом наименьших квадратов на интервале \(10^3\le n\le 10^5\). Этот инструмент уменьшает максимальную ошибку на интервале калибровки до \(\approx 0.006\%\). Заключительный поиск по окрестности гарантирует точную идентификацию простого числа. Константы \(A\) и \(B\) не выводятся из YUCT и должны рассматриваться как чисто эмпирические; их применимость вне калибровочного диапазона не гарантирована.
	
	\subsection{Остаточные осцилляции и нули \(\zeta(s)\)}
	
	Разность между истинными простыми числами и теоретической формулой YUCT состоит из детерминированных осцилляций, которые почти идеально (\(R^2>0.99\)) описываются вкладами нетривиальных нулей дзета-функции Римана. Это наблюдение, о котором сообщается в Приложении PrimeN, указывает на то, что эти нули являются собственными частотами поля координации. Полное беcпараметрическое замкнутое выражение для \(p_n\) без какого-либо поиска потребовало бы вывода нулей на основе YUCT и остаётся для будущей работы.
	
	\textbf{Статус результатов:}
	\begin{itemize}
		\item Теоретическая YUCT-формула: \textbf{Теорема} (асимптотическая).
		\item Эмпирическая YUCT-формула: \textbf{Эмпирический инструмент}, не теорема.
		\item Спектральная связь с нулями \(\zeta(s)\): \textbf{Эмпирическое наблюдение}, теоретическое объяснение открыто.
	\end{itemize}
	
	%%%%%%%%%%%%%%%%%%%%%%%%%%%%%%%
	% ЧАСТЬ IV: 19-МЕРНАЯ РОДИТЕЛЬСКАЯ ТЕОРИЯ
	%%%%%%%%%%%%%%%%%%%%%%%%%%%%%%%
	\part{19-мерная родительская теория}
	
	\section{Явный вид 19-мерного действия}
	\label{sec:19Dlagrangian}
	\[
	S_{19} = \frac{1}{2\kappa_{19}} \int d^{19}X \sqrt{-G}\; R(G) \;+\; \int d^{19}X \sqrt{-G}\; \mathcal{L}_{\text{coord}},
	\]
	с
	\[
	\mathcal{L}_{\text{coord}} = -\frac{1}{4} \Tr(\Psi_{MN}\Psi^{MN}) + \frac{1}{2} G^{MN}\,\partial_M K_{\mathrm{eff}}\,\partial_N K_{\mathrm{eff}} - V(K_{\mathrm{eff}}) + \dots
	\]
	\textbf{Статус:} постулат.
	
	%%%%%%%%%%%%%%%%%%%%%%%%%%%%%%%
	% ЧАСТЬ V: ПРОВЕРЯЕМЫЕ ПРЕДСКАЗАНИЯ И ПРОГРАММА
	%%%%%%%%%%%%%%%%%%%%%%%%%%%%%%%
	\part{Проверяемые предсказания и дальнейшая работа}
	
	\section{Конкретные экспериментальные проверки}
	\begin{itemize}
		\item \textbf{Тензор-скалярное отношение:} \(r \approx 0.07\).
		\item \textbf{Субмикрометровая гравитация:} \(\Delta g/g \sim 10^{-12}\) на 10~нм.
		\item \textbf{Неравенство CHSH:} нарушения не превышают \(2\sqrt{2}\).
		\item \textbf{Космический бюджет:} \(\Omega_\Lambda = 2/3\), \(\Omega_{\coord} = 0.283\), \(\Omega_b = 0.05\).
		\item \textbf{Температурная зависимость \(G\):} \(\Delta G/G \sim 10^{-16}\).
		\item \textbf{Асимптотическое масштабирование простых чисел:} локальный показатель \(\gamma\) должен сходиться к \(2/3\) при \(n\to\infty\) (уже проверено для \(n\le 5\times10^5\)).
	\end{itemize}
	
	\section{Программа полного вывода}
	\begin{enumerate}
		\item Построение внутреннего многообразия \(X_{15}\) и вычисление параметров масс фермионов.
		\item Вывод эффективного действия из 19-мерной теории.
		\item Ab initio статистическая механика словаря.
		\item Вывод нетривиальных нулей \(\zeta(s)\) из спектральных свойств поля координации.
	\end{enumerate}
	
	\section{Заключение}
	Мы представили самосогласованный вывод универсального показателя \(\beta = 2/3\) из аксиом YUCT, включая явный анализ функции стоимости и теорему об оптимальном транспорте. Теперь показано, что тот же показатель определяет асимптотическое распределение простых чисел, что как укрепляет теорию, так и открывает новый мост между теорией чисел и физикой координации. Остальные элементы теории четко обозначены как гипотезы или эффективные описания, что создает прочную основу для дальнейших исследований.
	
	\begin{center}
		\textit{Вселенная — это самоскоординирующийся словарь.}
	\end{center}
	
	\begin{thebibliography}{99}
		\bibitem{YakushevLaw2026}
		А.\,В.\,Якушев, \emph{Закон координации Якушева: Фундаментальный закон реальности как фрактальная сеть координации}, Zenodo (2026), DOI:10.5281/zenodo.18444598.
		\bibitem{AppendixL}
		А.\,В.\,Якушев, ``Приложение L: Фрактальное масштабирование ошибок координации — универсальный закон от ДНК до космологии,'' в \emph{YUCT}, Zenodo (2026).
		\bibitem{AppendixY}
		А.\,В.\,Якушев, ``Приложение Y: Математические основы YUCT — вывод из первых принципов,'' в \emph{YUCT}, Zenodo (2026).
		\bibitem{AppendixPrimeN}
		А.\,В.\,Якушев, ``Приложение PrimeN: Координационная лестница простых чисел YUCT,'' в \emph{YUCT}, Zenodo (2026).
		\bibitem{Rosser1941}
		J.\,B.\,Rosser, ``The \(n\)-th Prime,'' \emph{Amer. Math. Monthly} \textbf{48}, 607--611 (1941).
		\bibitem{West1997}
		West, G. B., Brown, J. H., \& Enquist, B. J. (1997). ``A general model for the origin of allometric scaling laws in biology.'' \emph{Science} \textbf{276}, 122--126.
		\bibitem{Banavar1999}
		Banavar, J. R., Maritan, A., \& Rinaldo, A. (1999). ``Size and form in efficient transportation networks.'' \emph{Nature} \textbf{399}, 130--132.
	\end{thebibliography}
	
\end{document}